\subsection{Decision Model과 Cross-Encoder의 재정렬 성능·지연시간 비교 (RQ2)}

\label{sec:ch4_3}

Jev는 재정렬 후 상위 10개 후보(shortlist)의 gold Recall을 5.16점 높였고(CI [+2.61, +7.94]점; CE 69.47점, Jev 74.63점), 재정렬 단계 지연은 약 2.25배였다. 모든 gold가 shortlist에 남은 문항도 112개(49.56\%)에서 130개(57.52\%)로 7.96\%p 증가하였다(CI [+3.89, +12.38]\%p). 두 구간은 결과 관찰 후 같은 절차로 추가 산출하였다(S1).

이 개선이 최종 완전 회수로 모두 이어지지는 않았다. 목록당 하나만 선택하는 제약을 반영하면 shortlist에서 모든 gold를 동시에 선택할 수 있는 문항은 CE 104개, Jev 120개였으나 실제 완전 회수는 77개와 79개였다. shortlist에 남은 gold 출현 중 최종 선택된 비율도 CE 79.52\%(233/293), Jev 77.64\%(250/322)로 비슷하여, Jev의 이득은 주로 재정렬 단계의 후보 유지에서 나타났고 선택 단계의 추가 이득은 관찰되지 않았다.

지연은 문항 내 중복을 제거한 질문–후보 쌍의 점수화부터 목록별 상위 10개 구성까지로 정의하고, 캐시 없이 3회 반복한 문항별 평균을 분석 단위로 하였다(측정 절차는 S9). CE는 로컬 MPS(float32·배치 50), Jev는 원격 API(동시 요청 4개)로 실행하였다.

\begin{table}[H]
\centering
\caption{재정렬 단계 지연(60문항, 초). 각 문항의 3회 평균 분포에서 계산하였다.}
\label{tab:ch4_3}
\small
\setlength{\tabcolsep}{3pt}
\begin{tabular}{@{}lrr@{}}
\hline
통계량 & CE & Jev \\
\hline
평균 & 5.01 & 10.65 \\
중앙값 & 4.82 & 10.79 \\
문항 평균의 제95백분위수 & 8.87 & 16.53 \\
\hline
\end{tabular}
\end{table}

Jev는 문항당 평균 5.64초 길었고(95\% CI [+5.15, +6.17]초), 문항별 Jev/CE 비율의 기하평균은 2.25배(CI [2.11, 2.39])였다. 이 비교는 로컬 CE와 원격 Jev라는 배포 조건에 의존하므로 모델 고유 속도로 일반화할 수 없고, 예상 법령 생성·BM25·LLM 선택·답변 생성·모델 로드를 제외하므로 전체 응답시간의 차이를 뜻하지 않는다.

반복 실행에서 기존 품질 실험과 모든 목록의 상위 10개가 순서까지 일치한 비율은 CE 100\%(180/180회), Jev 1.11\%(2/180회)였고, Jev의 순서 무관 후보 겹침은 평균 91.51\%였다(정의와 분포는 S9). 따라서 재호출 시 동일한 선택·답변 품질은 보장되지 않는다.
